\documentclass[11pt]{article}
\usepackage{fontspec}
\usepackage[french,english]{babel}
\usepackage{amsmath}
\usepackage{amssymb}
\usepackage{amsthm}
\usepackage{mathrsfs}
\usepackage{tikz-cd}
\usepackage{enumitem}

\newcommand{\printedpage}[1]{\marginpar{\footnotesize p.~#1}}
\newcommand{\msp}[1]{\marginpar{\footnotesize [#1]}}
\newcommand{\limind}{\mathop{\underrightarrow{\mathrm{lim}}}}
\newcommand{\limproj}{\mathop{\underleftarrow{\mathrm{lim}}}}
\newcommand{\M}{\mathcal{M}}
\newcommand{\Fl}{\mathrm{Fl}}
\newcommand{\Ob}{\mathrm{Ob}}
\newcommand{\Cof}{\mathrm{Cof}}
\newcommand{\Fib}{\mathrm{Fib}}
\newcommand{\univ}{\mathrm{univ}}
\newcommand{\HO}{\mathrm{HO}}

\begin{document}

\printedpage{24}
% head: A. Grothendieck : Dérivateurs, Ch. XIII, Édition des Universités de Montpellier II et Paris VII
to deduce a similar factorization for the arrows $f$ of $\M(I)$, as
\[
\begin{aligned}
f = pi \quad \text{with}\quad & i \in TC(I), && p \in F(I)\\
\text{or}\quad & i \in C(I), && p \in TF(I),
\end{aligned}
\]
where $C(I)$, $TC(I)$, $F(I)$, $TF(I)$ are here defined via $C$, $TC$, $F$, $TF$ ``argument by argument''. But the trouble is that by proceeding in this naive way, we will no longer necessarily have $C(I) \longleftrightarrow TF(I)$, $TC(I) \longleftrightarrow F(I)$. We have to shrink $TF(I)$ and $F(I)$ considerably, to have an RLP property with respect to $C(I)$ resp. $TC(I)$.

So I can hardly see any other way out, in this general context, than by assuming the stability of $W$ under filtering $\limind$ (which is after all not very nasty).

\noindent\msp{page 42}

Willy-nilly, we come back to the situation considered in paragraph 2. We are going to restate the essentials of what was sensed then:

\textbf{Theorem 2} (still a little conjectural). {\itshape Let $(\M,W)$ be a category of models, with $W$ (as always) weakly saturated, but in addition we assume $\M$ stable under small $\limind$ and under finite $\limproj$, and $W$ stable under filtering $\limind$, finally $\M$ accessible. Consider the following conditions:

\begin{enumerate}
\item[\textup{a)}] There exists a weak Quillen structure for $W$, left exact, right closed, $(C,TC,F,TF)$.
\item[\textup{a$'$)}] As in \textup{a)}, but with in addition $C$, $TC$ left closed (\textup{i.e.} we have a closed weak Quillen structure), and $C$, $TC$ of the form $\widetilde{C_0}$, $\widetilde{TC_0}$, where $C_0$ and $TC_0$ are small.
\item[\textup{a$''$)}] There exist $C_0 \subset \Cof_W$, $TC_0 \subset W^{\univ}$, small subsets, such that $C_{0*} \subset W$, $(TC_0)_* \subset \Fib_W$.
\item[\textup{b)}] We have
\[
(\Cof_W)_* \;\subset\; W_{\univ}, \qquad\quad (W^{\univ})_* \;\subset\; \Fib_W
\]
(whence it follows, $\Cof_W$ and $W^{\univ}$ being left closed, that
\[
(W_{\univ})^* \subset \Cof_W, \qquad\quad (\Fib_W)^* \subset W^{\univ} \qquad\quad \text{\itshape ).}
\]
\end{enumerate}

\noindent\msp{page 43}

These conditions imply that for every $I$ in \textup{Cat}, there exists a left exact and closed weak Quillen structure for $W(I)$ in $\M(I)$. Such a structure is defined, in terms of a similar structure $C$, $TC = C \cap W$, $F = (TC)_*$, $TF = C_*$ for $W$, by}
\[
C(I) \;=\; \Bigl\{ f \in \Fl(\M(I)) \Bigm| \alpha_s^*(f) \in C \ \ \forall\, s \in \Ob I \Bigr\},
\]

\printedpage{25}
% head: A. Grothendieck : Dérivateurs, Ch. XIII, Édition des Universités de Montpellier II et Paris VII
\noindent\textit{then}
\[
\begin{aligned}
TC(I) \;&=\; C(I)\cap W(I) \;=\; \Bigl\{ f \in \Fl(\M(I)) \Bigm| \alpha_s^*(f) \in TC \ \ \forall\, s \in \Ob I \Bigr\}\\
F(I) \;&=\; (TC(I))_*\\
TF(I) \;&=\; C(I)_*\,.
\end{aligned}
\]

\textbf{Corollary-Scholie.} \textit{These conditions imply that for every $I$ in} Cat, \textit{$(\M(I),W(I))$ gives rise to a good theory of $W(I)$-homotopic fibers and cofibers, with $W(I)$-cartesian $W(I)$-exact squares, and $W(I)$-cocartesian $W(I)$-exact squares – hence with long} exact \textit{sequences of cosuspension and of suspension, in the case where $\M$, hence also $\M(I)$, is a category with a zero object.}

\noindent\msp{page 44}

I will take care of the proof of the preceding theorem, as well as of the factorization theorem of paragraph 1, when I have my notes on accessible categories and functors and cardinal filtrations. In the meantime, I already feel like making a little more explicit the formalism one arrives at by means of the hypotheses made above. (I am pretty sure at least that condition $a''$) suffices to imply all the others. And it is indeed in the form $a''$) that these conditions are going to be verified in practice, I presume, notably in the case where $\M = $ Cat.)

\section*{4\quad The variances $u^*$, $u_*$, $u_{!}$ for $(\M(I),W(I))$ (I variable)}

I assume the conditions of the preceding paragraph are satisfied for $(\M,W)$ (cf. theorem p. 42), and choose a closed and left exact weak Quillen structure $(C,TC,F,TF)$ for $W$, whence for every $I$ in Cat

\noindent\msp{page 45}

\noindent a closed and left exact weak Quillen structure $(C(I),TC(I),F(I),TF(I))$ for $W(I)$, in $\M(I)$, described on page 43. Now let
\[
u \;:\; I' \longrightarrow I
\]
be an arrow in Cat, whence an arrow
\[
u^* \;:\; (\M(I),W(I)) \;\longrightarrow\; (\M(I'),W(I')).
\]
We know, since $\M$ is stable under arbitrary $\limind$ and also (being moreover accessible) under arbitrary $\limproj$, that $u^*$ admits left and right adjoint functors,
\[
u_{!},\ u_* \;:\; \M(I') \;\longrightarrow\; \M(I),
\]
\printedpage{26}
% head: A. Grothendieck : Dérivateurs, Ch. XIII, Édition des Universités de Montpellier II et Paris VII
and I propose here to summarize the most immediate properties of $u^*$ and of $u_{!}$, $u_*$ with respect to the weak Quillen structures under consideration.

\textbf{Proposition 1:} \textit{With the preceding notations, we have the relations:}
\[
u^*(C(I)) \subset C(I'), \qquad u^*(TC(I)) \subset TC(I'). \tag*{(1)}
\]
\[
u_*(F(I')) \subset F(I), \qquad u_*(TF(I')) \subset TF(I). \tag*{(2)}
\]
\[
\begin{aligned}
u^*(F(I)) \subset F(I') &\iff u_{!}(TC(I')) \subset TC(I)\\
u^*(TF(I)) \subset TF(I') &\iff u_{!}(C(I')) \subset C(I).
\end{aligned} \tag*{(3)}
\]

On the other hand, I do not know anything positive to say about the inclusions
\[
\begin{aligned}
u_*(C(I')) \;&\subset\; C(I), \;\; u_*(TC(I')) \;\subset\; TC(I)\\
u_{!}(F(I')) \;&\subset\; F(I), \;\; u_{!}(TF(I')) \;\subset\; TF(I).
\end{aligned}
\]
Except that for a $u$ that is at all general ($u$ not an isomorphism, say), $u_*$ has no desire

\noindent\msp{page 46}

\noindent to be compatible with the structures $C$, $TC$ (a property of homotopical right exactness, morally, since $u_*$ is left exact and in general not at all right exact), nor $u_{!}$ with the structures $F$, $TF$ (dual comment to be given, $u_{!}$ is \textit{right} exact, and does not feel like being ``homotopically left exact''). % [sic] exacte -> exact, accord avec « foncteur » / « exact à gauche »

\textsc{Proof} of (1), (2), (3). Relations (1) are obvious by definition. Relations (2), (3) are scarcely less so, in view of the following general sorital property, whose place would be in paragraph 1:

\textbf{Proposition 2.} \textit{Let}
\[
\begin{tikzcd}
C' \arrow[r, shift left, "u_{!}"] & C \arrow[l, shift left, "u^*"]
\end{tikzcd}
\]
\textit{be a pair of adjoint functors $(u_{!},u^*)$, and let $\Phi' \subset \Fl(C')$, $\Psi \subset \Fl(C)$. Then}
\[
(u_{!}(\Phi') \longleftrightarrow \Psi) \qquad\quad \iff \qquad\quad (\Phi' \longleftrightarrow u^*(\Psi)).
\]

\textbf{Corollary.} (This is the same statement, with only the notation changed.) \textit{Suppose that $u^* : C \longrightarrow C'$ admits a} right \textit{adjoint $u_*$ (not necessarily a left adjoint)}
\[
\begin{tikzcd}
C' \arrow[r, shift right, "u_*"'] & C, \arrow[l, shift right, "u^*"']
\end{tikzcd}
\]

\printedpage{27}
% head: A. Grothendieck : Dérivateurs, Ch. XIII, Édition des Universités de Montpellier II et Paris VII
\noindent\textit{let $\Phi \subset \Fl(C)$, $\Psi' \subset \Fl(C')$, then}
\[
(u^*(\Phi) \longleftrightarrow \Psi') \qquad\quad \iff \qquad\quad (\Phi \longleftrightarrow u_*(\Psi')).
\]

This is proposition 2 for \begin{tikzcd} C'_0 \arrow[r, shift left, "u_{0\,!}"] & C_0 \arrow[l, shift left, "u_0^*"] \end{tikzcd}, with $C'_0 = C$, $C_0 = C'$, $(u_{0\,!} : C \longrightarrow C') = u^*$, and $(u_0^* : C' \longrightarrow C) = u_*$.

\noindent\msp{page 47}

For the inclusions (2), we use the corollary of proposition 2 for the pairs $\Phi = TC(I)$, $\Psi' = F(I')$, resp. $\Phi = C(I)$, $\Psi' = F(I')$ \texttt{[rather} $\Psi' = TF(I')$\texttt{]}, and apply (1) which gives $u^*(\Phi) \longleftrightarrow \Psi'$, whence $\Phi \longleftrightarrow u_*(\Psi')$, which is nothing other than (2). For the equivalences (3), we use proposition 2, putting $\Psi = F(I)$, $\Phi' = TC(I')$ resp. $\Psi = TF(I)$, $\Phi' = C(I')$.

\textbf{Corollary 1} (of proposition 1): \textit{Suppose $u : I' \longrightarrow I$ is such that for every $i \in \Ob I$, the category $i\backslash I'$ is discrete} (\footnote[27]{More generally, it suffices that every connected component $Z_\alpha$ of $i\backslash I'$ have an initial object $e_\alpha$, so that $\limind\limits_{i\backslash I'} X'(i') \simeq \coprod_{\alpha\in\pi_0(i\backslash I')} X'(e_\alpha)$. This is the case in particular if $I' \longrightarrow I$ is a local isomorphism, e.g. an open immersion.}) \textit{(e.g. $I'$ discrete). Then we have}
\[
\left\{
\begin{aligned}
u_{!}(C(I')) \;&\subset\; C(I), \;\; u_{!}(TC(I')) \;\subset\; TC(I)\\
u^*(F(I)) \;&\subset\; F(I'), \;\; u^*(TF(I)) \;\subset\; TF(I').
\end{aligned}
\right.
\]

Indeed, the relations on the second line are equivalent (because of (3)) to those on the first. We have to prove that if $f$ in $\Fl\M(I')$ is in $C(I')$ resp. $TC(I')$, $u_{!}(f)$ is in $C(I)$ resp. $TC(I)$, which also means that $\alpha_s^*(u_{!}(f))$ is in $C$ resp. $TC$, for every $s \in \Ob I$. Now the functor
\[
\alpha_s^* u_{!} \;:\; \M(I') \;\longrightarrow\; \M
\]
is given explicitly by the well-known formula
\[
\alpha_s^*(u_{!}(X')) \;\overset{\text{def}}{=}\; u_{!}(X')(s) \;\simeq\; \limind\limits_{i'\in s\backslash I'} X'(i').
\]

Since

\noindent\msp{page 48}

\noindent $s\backslash I'$ is discrete, we therefore have the isomorphism functorial in $X'$ in $\M(I')$
\[
\alpha_s^*(u_{!}(X')) \;\simeq \coprod_{(s',\, x \,:\, s \,\longrightarrow\, u(s')) \text{ in } \Ob\,(s\backslash I')} X'(s'),
\]
so if $f \,:\, X' \longrightarrow Y'$, $\alpha_s^*(u_{!}(f))$ is made explicit as a direct sum in $\underline{\mathrm{Fl}}(\M)$ of arrows of the form
\[
f_{s'} \;:\; X'(s') \;\longrightarrow\; Y'(s'),
\]
\printedpage{28}
% head: A. Grothendieck : Dérivateurs, Ch. XIII, Édition des Universités de Montpellier II et Paris VII
which are in $C$ resp. $TC$ since $f$ is in $C(I')$ resp. $TC(I')$. As $C$, $TC$ are stable under arbitrary direct sums ($C$, $TC$ being left closed), we win.

\textbf{Corollary 2.} \textit{Let $I$ be in} Cat, \textit{and $f \in \Fl(\M(I))$. For $f$ to be in $F(I)$} (resp. \textit{in $TF(I)$\,), it} is necessary \textit{that $\forall\, s \in \Ob I$, $f_s : X(s) \longrightarrow Y(s)$ be so.} (But this condition does not look sufficient at all!) \textit{In particular, let $X \in \Ob\M(I)$, then $X \in \Ob\M(I)_f \Longrightarrow X(s) \in \Ob\M_f \ \ \forall\, s \in \Ob I$.}

We apply corollary 1 to $\alpha_s : e \longrightarrow I$.

\textbf{Scholie.} Proposition 1 (relation (1)) shows that the functor $u^*$ transforms $W$-cocartesian squares into likewise, long exact sequences of suspension into likewise – this is ``the homotopical right exactness'' of $u^*$. We expect the same to hold for cartesian squares and long exact sequences

\noindent\msp{page 49}

\noindent of cosuspension, but to prove it some work will be needed, which we will do in what follows.

Relations (2) show that the functor $u_*$ transforms $W$-cartesian squares into likewise, hence long exact sequences of cosuspension into likewise. I do \textit{not} expect the same to hold in general for $W$-cocartesian squares – \textit{i.e.} that $u_*$ be ``homotopically right exact'', and not only left exact. But I should at least do some examples (with $I'$, $I$ finite ordered categories).

As for $u_{!}$, the proposition says nothing very tangible about it. Nevertheless we expect $u_{!}$ to transform $W$-cocartesian squares into likewise – that is the least one can expect! Or perhaps I am mistaken, in confusing the functors $u_*$, $u_{!}$ considered here, with the similar functors
\[
\begin{tikzcd}
\HO_W(I') \arrow[r, shift left, "u_{!}^W"] \arrow[r, shift right, "u_*^W"'] & \HO_W(I),
\end{tikzcd}
\]
left resp. right adjoints to
\[
\begin{tikzcd}
\HO_W(I') & \HO_W(I), \arrow[l, "u_W^*"']
\end{tikzcd}
\]
induced by
\[
u^* \;:\; (\M(I),W(I)) \;\longrightarrow\; (\M(I'),W(I')).
\]
But in general $u_*$, $u_{!}$ are not compatible

\noindent\msp{page 50}

\noindent with the localizers $W(I')$, $W(I)$, so they must not be confused with the functors $u_{!}^W$, $u_*^W$. So perhaps I was hasty, in what I said about $u_*$ and $u_{!}$. I will be obliged to study the situation very closely, when I try to construct the functors $u_*^W$ and $u_{!}^W$. For the moment, I still have no assurance that the conditions considered for $(\M,W)$ suffice to ensure the existence of $u_*^W$, $u_{!}^W$ – not even for $I$, $I'$ finite ordered sets!

\textbf{Corollary 1.} \textit{Let $X \in \M(I)_{\mathrm{f}}$, then for every $\alpha : s \longrightarrow t$ in $I$ which is} \texttt{[a]} monomorphism, \textit{$X(\alpha) : X(t) \longrightarrow X(s)$ is in $F$.}

% foot: Version 25 janvier 2004

\end{document}
